\documentclass[12pt, a4paper]{article}

% --- Pakiety ---
\usepackage[utf8]{inputenc}
\usepackage[T1]{fontenc}
\usepackage[polish]{babel}
\usepackage{amsmath, amssymb}
\usepackage{graphicx}
\usepackage{float}
\usepackage{geometry}
\usepackage{hyperref}
\usepackage{booktabs}
\usepackage{caption}
\usepackage{listings}
\usepackage{xcolor}
\usepackage{parskip}

\geometry{
  a4paper,
  left=2.5cm,
  right=2.5cm,
  top=2.5cm,
  bottom=2.5cm
}

% --- Zakaz dzielenia wyrazów ---
\hyphenpenalty=10000
\exhyphenpenalty=10000
\emergencystretch=3em

\hypersetup{
  colorlinks=true,
  linkcolor=blue,
  urlcolor=blue,
  citecolor=blue
}

% --- Zmiana nazwy float figure: Rysunek -> Wykres ---
\addto\captionspolish{\renewcommand{\figurename}{Wykres}}

% --- Styl listingów ---
\lstset{
  language=Python,
  basicstyle=\ttfamily\small,
  keywordstyle=\color{blue}\bfseries,
  commentstyle=\color{gray}\itshape,
  stringstyle=\color{orange!80!black},
  numbers=left,
  numberstyle=\tiny\color{gray},
  numbersep=8pt,
  frame=single,
  breaklines=true,
  breakatwhitespace=false,
  captionpos=b,
  showstringspaces=false,
  inputencoding=utf8,
  extendedchars=true,
  literate={ą}{{\k{a}}}1 {ę}{{\k{e}}}1 {ó}{{\'o}}1
            {ś}{{\'s}}1 {ł}{{\l{}}}1 {ź}{{\'z}}1
            {ż}{{\.z}}1 {ć}{{\'c}}1 {ń}{{\'n}}1
            {Ą}{{\k{A}}}1 {Ę}{{\k{E}}}1 {Ó}{{\'O}}1
            {Ś}{{\'S}}1 {Ł}{{L}}1,
}

% --- Strona tytułowa ---
\title{
  \textbf{Metody Obliczeniowe w Nauce i Technice} \\[0.5em]
  \large Laboratorium 3 \textemdash{} Interpolacja
}
\author{Mateusz Stoch}
\date{1 kwietnia 2026}

\begin{document}

\maketitle

% ============================================================
\section{Zadanie 1 \textemdash{} Interpolacja populacji USA}
% ============================================================

% --- 1.1 Opis zadania ---
\subsection{Opis zadania}

Zadanie polega na wyznaczeniu wielomianu interpolacyjnego ósmego stopnia
na podstawie danych o~populacji USA z~lat 1900--1980 (9 punktów, co 10 lat).
Rozważane są cztery różne bazy funkcji bazowych $\phi_j(t)$:
\begin{align}
  \phi_j(t) &= t^{j-1}, \label{eq:baza1} \\
  \phi_j(t) &= (t - 1900)^{j-1}, \label{eq:baza2} \\
  \phi_j(t) &= (t - 1940)^{j-1}, \label{eq:baza3} \\
  \phi_j(t) &= \left(\frac{t - 1940}{40}\right)^{j-1}. \label{eq:baza4}
\end{align}
Na podstawie najlepiej uwarunkowanej bazy wyznacza się współczynniki wielomianu
i~dokonuje ekstrapolacji do roku 1990. Zadanie obejmuje też wyznaczenie
wielomianów Lagrange'a i~Newtona, a~następnie analizę wpływu zaokrąglenia
danych na jakość interpolacji.

% --- 1.2 Argumentacja ---
\subsection{Argumentacja rozwiązania}

\subsubsection{Macierz Vandermonde'a i współczynnik uwarunkowania}

Dla każdej bazy budowana jest macierz Vandermonde'a $\mathbf{V}$, gdzie
$V_{ij} = \phi_j(t_i)$. Rozwiązując układ $\mathbf{V}\mathbf{c} = \mathbf{y}$,
dostajemy współczynniki wielomianu. Problem w~tym, że dla bazy~(\ref{eq:baza1})
wartości $t_i$ są rzędu 1900--1980, więc potęgując je do 8 dostajemy liczby
rzędu $10^{26}$~--- macierz jest skrajnie źle uwarunkowana.

Przesuwanie i~skalowanie argumentu (bazy 2--4) drastycznie poprawia
uwarunkowanie. Baza~(\ref{eq:baza4}) jest najlepsza, bo skaluje wartości do
przedziału $[-1, 1]$.
\newpage
\begin{lstlisting}[caption={Budowa macierzy Vandermonde'a}]
t = years.astype(float)
V1 = np.vander(t,              N=9, increasing=True)
V2 = np.vander(t - 1900,       N=9, increasing=True)
V3 = np.vander(t - 1940,       N=9, increasing=True)
V4 = np.vander((t-1940)/40,    N=9, increasing=True)
\end{lstlisting}

\subsubsection{Schemat Hornera}

Do obliczania wartości wielomianu w~wielu punktach używa się schematu Hornera,
który jest numerycznie stabilniejszy i~szybszy niż bezpośrednie sumowanie potęg:
\begin{equation}
  p(x) = c_0 + x\bigl(c_1 + x(c_2 + \cdots + x\,c_n)\bigr).
\end{equation}

\begin{lstlisting}[caption={Schemat Hornera dla bazy 4}]
def horner_basis4(t_eval, c, center=1940, scale=40):
    x = (t_eval - center) / scale
    result = np.zeros_like(t_eval, dtype=float)
    for coeff in reversed(c):
        result = result * x + coeff
    return result
\end{lstlisting}

\subsubsection{Wielomian Lagrange'a}

Wzór Lagrange'a wyraża wielomian jako kombinację liniową funkcji bazowych $L_i$:
\begin{equation}
  p(t) = \sum_{i=0}^{n} y_i\, L_i(t), \qquad
  L_i(t) = \prod_{j \neq i} \frac{t - t_j}{t_i - t_j}.
\end{equation}
Nie wymaga rozwiązywania układu równań, ale jest wolniejszy dla dużej liczby węzłów.

\subsubsection{Wielomian Newtona (iloraz różnicowy)}

Wielomian Newtona jest zbudowany na ilorazach różnicowych:
\begin{equation}
  p(t) = c_0 + c_1(t{-}t_0) + c_2(t{-}t_0)(t{-}t_1) + \cdots
\end{equation}
Zaletą tej postaci jest to, że dodanie nowego węzła wymaga tylko jednej dodatkowej
operacji~--- nie trzeba przebudowywać całego wielomianu.

% --- 1.3 Wykresy i wyniki ---
\subsection{Wykresy i wyniki liczbowe}

\begin{table}[H]
  \centering
  \caption{Współczynniki uwarunkowania macierzy Vandermonde'a dla czterech baz.}
  \begin{tabular}{lll}
    \toprule
    Baza & Definicja & $\operatorname{cond}(\mathbf{V})$ \\
    \midrule
    Baza 1 & $t^{j-1}$                  & $\approx 5{,}36 \times 10^{36}$ \\
    Baza 2 & $(t-1900)^{j-1}$           & $\approx 5{,}87 \times 10^{15}$ \\
    Baza 3 & $(t-1940)^{j-1}$           & $\approx 9{,}32 \times 10^{12}$ \\
    Baza 4 & $((t-1940)/40)^{j-1}$      & $\approx 1{,}61 \times 10^{3}$  \\
    \bottomrule
  \end{tabular}
  \label{tab:cond_van}
\end{table}

Najlepiej uwarunkowana jest Baza 4~--- to ona została użyta do wyznaczenia
współczynników i~ekstrapolacji.

\begin{figure}[H]
  \centering
  \includegraphics[width=0.85\textwidth]{wykres1.png}
  \caption{Wielomian interpolacyjny (baza 4) na przedziale 1900--1990
           z~zaznaczonymi węzłami interpolacji.}
  \label{fig:interpolacja_baza4}
\end{figure}

\begin{figure}[H]
  \centering
  \includegraphics[width=0.85\textwidth]{wykres2.png}
  \caption{Wielomian interpolacyjny Lagrange'a na tych samych danych.}
  \label{fig:lagrange1}
\end{figure}

\begin{figure}[H]
  \centering
  \includegraphics[width=0.85\textwidth]{wykres3.png}
  \caption{Wielomian interpolacyjny Newtona.}
  \label{fig:newton}
\end{figure}

\begin{table}[H]
  \centering
  \caption{Ekstrapolacja do roku 1990 -- porównanie wyników.}
  \begin{tabular}{lrr}
    \toprule
    Dane wejściowe & Predykcja (1990) & Błąd względny \\
    \midrule
    Oryginalne            & $82\,749\,141$  & $\approx 66{,}73\%$ \\
    Zaokrąglone do 1 mln  & $109\,000\,000$ & $\approx 56{,}17\%$   \\
    Wartość prawdziwa     & $248\,709\,873$ & ---                 \\
    \bottomrule
  \end{tabular}
  \label{tab:ekstrapolacja}
\end{table}

\noindent
Jak widać, wielomian interpolacyjny stopnia~8 bardzo dobrze odwzorowuje
dane w~przedziale węzłów (1900--1980), natomiast ekstrapolacja do~roku 1990
daje znacznie błędny wynik (błąd $\approx 66{,}7\%$)~--- typowa słabość
wielomianów wysokich stopni poza zakresem interpolacji. Zaokrąglenie danych
do~miliona zmienia współczynniki wielomianu, co jeszcze bardziej wpływa
na~wynik ekstrapolacji.


% --- 1.4 Wnioski ---
\subsection{Wnioski}

\begin{enumerate}
  \item Wybór bazy ma ogromny wpływ na uwarunkowanie macierzy Vandermonde'a.
        Baza~1 (surowe lata) jest praktycznie bezużyteczna numerycznie
        (cond $\approx 5{,}36 \times 10^{36}$, patrz Tabela~\ref{tab:cond_van}).
        Baza~4 z~przesunięciem i~skalowaniem daje uwarunkowanie rzędu $10^3$,
        co przekłada się na stabilne obliczenia.

  \item Wszystkie trzy metody (baza~4, Lagrange, Newton) dają ten sam wielomian
        stopnia~8 i~bardzo dobrze odwzorowują dane w~przedziale węzłów
        (1900--1980). Różnice w~wynikach wynikają jedynie z~błędów zaokrąglenia.

  \item Ekstrapolacja jest zawodna~--- nawet przy dobrej bazie. Wielomian
        stopnia~8 dla danych oryginalnych przewidział populację w~1990 roku
        na poziomie $82{,}7$~mln, przy prawdziwej wartości $248{,}7$~mln
        (błąd $\approx 66{,}7\%$, patrz Tabela~\ref{tab:ekstrapolacja}).
        Poza zakresem węzłów wielomian zachowuje się nieprzewidywalnie.

  \item Zaokrąglenie danych do~1~mln zmieniło predykcję na $109$~mln
        (błąd $\approx 56{,}2\%$). Choć błąd jest mniejszy, to wynik dalej
        jest bardzo daleki od rzeczywistości. Wynika to z~wrażliwości
        współczynników wielomianu na perturbacje w~danych wejściowych.
\end{enumerate}

% ============================================================
\section{Zadanie 2 \textemdash{} Interpolacja produkcji cytrusów}
% ============================================================

% --- 2.1 Opis zadania ---
\subsection{Opis zadania}

Dane dotyczą rocznej produkcji cytrusów we Włoszech w~latach 1965--1991
(6 punktów, nierównomiernie rozłożonych). Należy oszacować produkcję
w~latach 1962, 1977 i~1992 (dwa przypadki interplanacji i~jeden ekstrapolacji)
używając:
\begin{itemize}
  \item kubicznych funkcji sklejanych (natural, not-a-knot, clamped),
  \item wielomianu interpolacyjnego Lagrange'a.
\end{itemize}
Prawdziwe wartości do porównania: 12\,380, 27\,403 i~32\,059
($\times 10^5$~kg).

% --- 2.2 Argumentacja ---
\subsection{Argumentacja rozwiązania}

\subsubsection{Kubiczne funkcje sklejane}

Zamiast jednego wielomianu wysokiego stopnia, funkcje sklejane (splines)
łączą kawałkowe wielomiany trzeciego stopnia tak, żeby na każdym węźle
były ciągłe wraz z~pierwszą i~drugą pochodną. Poprawia to zachowanie
na całym przedziale i~eliminuje efekt Rungego.

Różnią się warunkami brzegowymi (co dzieje się na końcach przedziału):
\begin{itemize}
  \item \textbf{Natural}: druga pochodna na końcach równa zero
        ($f''(x_0) = f''(x_n) = 0$). Spline „wygładza się" na końcach.
  \item \textbf{Not-a-knot}: trzecia pochodna jest ciągła na pierwszym
        i~ostatnim węźle wewnętrznym. Często daje dobre wyniki.
  \item \textbf{Clamped}: narzucamy wartość pierwszej pochodnej na końcach
        (tu: zero). Przydatne, gdy znamy nachylenie na krańcach.
\end{itemize}

\begin{lstlisting}[caption={Kubiczne funkcje sklejane -- trzy warianty}]
from scipy.interpolate import CubicSpline

cs_natural = CubicSpline(years2, production, bc_type='natural')
cs_nak     = CubicSpline(years2, production, bc_type='not-a-knot')
cs_clamped = CubicSpline(years2, production,
                         bc_type=((1, 0.0), (1, 0.0)))
\end{lstlisting}

\subsubsection{Wielomian Lagrange'a na cytrusach}

Użyto tej samej funkcji co w~zadaniu 1. Przy 6 węzłach otrzymujemy
wielomian piątego stopnia. W~przypadku ekstrapolacji (rok 1992 jest
poza zakresem danych) wielomian może mocno odbiegać od rzeczywistości.

% --- 2.3 Wykresy i wyniki ---
\subsection{Wykresy i wyniki liczbowe}

\begin{figure}[H]
  \centering
  \includegraphics[width=0.9\textwidth]{wykres4.png}
  \caption{Kubiczne funkcje sklejane (trzy warianty) na tle węzłów
           i~prawdziwych wartości dla lat 1962, 1977, 1992.}
  \label{fig:splines}
\end{figure}

\begin{figure}[H]
  \centering
  \includegraphics[width=0.9\textwidth]{wykres5.png}
  \caption{Wielomian interpolacyjny Lagrange'a dla danych o~cytrusach.}
  \label{fig:lagrange2}
\end{figure}

\begin{table}[H]
  \centering
  \caption{Wyestymowane wartości produkcji i~błędy względne dla roku 1962
           (ekstrapolacja wsteczna).}
  \begin{tabular}{lrr}
    \toprule
    Metoda & Estymacja [$10^5$ kg] & Błąd względny \\
    \midrule
    Wartość prawdziwa & 12\,380     & --- \\
    Natural spline    & $\approx 13\,285$ & $\approx 7{,}31\%$ \\
    Not-a-knot        & $\approx 5\,146$  & $\approx 58{,}43\%$ \\
    Clamped           & $\approx 24\,313$ & $\approx 96{,}39\%$ \\
    Lagrange          & $\approx -77\,765$ & $\approx 728{,}15\%$ \\
    \bottomrule
  \end{tabular}
  \label{tab:cytrusy_1962}
\end{table}

\begin{table}[H]
  \centering
  \caption{Wyestymowane wartości produkcji i~błędy względne dla roku 1977
           (interpolacja wewnątrz przedziału).}
  \begin{tabular}{lrr}
    \toprule
    Metoda & Estymacja [$10^5$ kg] & Błąd względny \\
    \midrule
    Wartość prawdziwa & 27\,403    & --- \\
    Natural spline    & $\approx 22\,934$ & $\approx 16{,}31\%$ \\
    Not-a-knot        & $\approx 22\,642$ & $\approx 17{,}37\%$ \\
    Clamped           & $\approx 23\,126$ & $\approx 15{,}61\%$ \\
    Lagrange          & $\approx 15\,405$ & $\approx 43{,}78\%$ \\
    \bottomrule
  \end{tabular}
  \label{tab:cytrusy_1977}
\end{table}

\begin{table}[H]
  \centering
  \caption{Wyestymowane wartości produkcji i~błędy względne dla roku 1992
           (ekstrapolacja do przodu).}
  \begin{tabular}{lrr}
    \toprule
    Metoda & Estymacja [$10^5$ kg] & Błąd względny \\
    \midrule
    Wartość prawdziwa & 32\,059    & --- \\
    Natural spline    & $\approx 37\,798$ & $\approx 17{,}90\%$ \\
    Not-a-knot        & $\approx 41\,894$ & $\approx 30{,}68\%$ \\
    Clamped           & $\approx 22\,166$ & $\approx 30{,}86\%$ \\
    Lagrange          & $\approx 43\,127$ & $\approx 34{,}52\%$ \\
    \bottomrule
  \end{tabular}
  \label{tab:cytrusy_1992}
\end{table}

% --- 2.4 Wnioski ---
\subsection{Wnioski}

\begin{enumerate}
  \item Spline naturalny wypadł najlepiej przy ekstrapolacji wstecznej
        (rok 1962, błąd $\approx 7{,}31\%$). Ma sensowne zachowanie na
        krańcach, bo zeruje drugą pochodną i~nie ``szaleje'' poza zakresem.
        Not-a-knot i~Clamped poradziły sobie znacznie gorzej
        (odpowiednio $\approx 58\%$ i~$\approx 96\%$ błędu).

  \item Przy interpolacji w~środku (rok 1977) wszystkie trzy spline'y
        dały zbliżone błędy (ok.~$15$--$17\%$). Wielomian Lagrange'a
        wypadł gorzej ($\approx 44\%$).

  \item Ekstrapolacja do przodu (rok 1992) to najtrudniejszy przypadek.
        Spline naturalny wypadł tu relatywnie najlepiej ($\approx 18\%$).
        Not-a-knot i~Clamped osiągnęły podobny błąd ($\approx 31\%$).
        Wielomian Lagrange'a dał $\approx 35\%$ błędu.

  \item Lagrange przy ekstrapolacji wstecznej (rok 1962) całkowicie
        zawiódł~--- estymacja wyniosła $-77\,765$, czyli wartość
        ujemną przy prawdziwej wartości $12\,380$. Błąd to
        ponad $728\%$. To pokazuje, jak niebezpieczny jest wielomian
        wysokiego stopnia poza zakresem danych.

  \item Funkcje sklejane są bezpieczniejsze od wielomianów wysokiego stopnia,
        bo działają lokalnie~--- zaburzenie w~jednym miejscu nie propaguje
        się na cały przedział.
\end{enumerate}

% ============================================================
\section{Podsumowanie}
% ============================================================

\begin{enumerate}
  \item \textbf{Wybór bazy ma znaczenie.} Przy macierzy Vandermonde'a złe
        uwarunkowanie (baza $t^j$, cond~$\approx 5{,}36 \times 10^{36}$)
        dosłownie niszczy wyniki. Przesunięcie i~skalowanie argumentu
        (baza $((t-1940)/40)^j$, cond~$\approx 1{,}61 \times 10^{3}$)
        poprawiło uwarunkowanie o~ponad 33 rzędy wielkości.

  \item \textbf{Lagrange, Newton i równanie normalne to ten sam wielomian.}
        Każda z~metod daje tę samą krzywą, różnią się tylko sposobem
        reprezentacji i~stabilnością numeryczną.

  \item \textbf{Ekstrapolacja jest zawodna.}
        Wielomian stopnia~8 interpolujący populację USA (zadanie~1)
        przy ekstrapolacji do roku 1990 dał błąd $\approx 66{,}7\%$.
        Wielomian Lagrange'a dla cytrusów (zadanie~2) przy ekstrapolacji
        wstecznej (rok~1962) zwrócił wartość ujemną~--- ponad $728\%$ błędu.
        Wyniki potwierdzają, że poza zakresem węzłów rezultaty mogą być
        zupełnie nierealne.

  \item \textbf{Spline'y są lepszym wyborem w~praktyce.}
        Przy interpolacji na rozłożonych nierównomiernie danych, kubiczne
        funkcje sklejane zazwyczaj dają spokojniejszą i~bardziej
        wiarygodną krzywą niż wielomian interpolacyjny.
\end{enumerate}

% ============================================================
\section*{Bibliografia}
\addcontentsline{toc}{section}{Bibliografia}
% ============================================================

\begin{thebibliography}{9}

  \bibitem{lab03}
  \textit{Wprowadzenie do laboratorium}, materia\l{}y dost\k{e}pne
  na platformie Teams w~katalogu lab03.

\end{thebibliography}

\end{document}
